\section{Discussion}
\label{sec:discussion}

\paragraph{The explanation was faithful but wrong in size and in kind}
The world model's explanation is faithful by construction: the baseline does not
move with the plan, the lift is zero when nothing is planned, and the two add up
to the forecast. In the known-answer study that explanation was right about what
each lever adds when plans were made at random, and wrong in two ways that no
check against the model could see. When merchants plan into strong weeks, the
lifts are a third too large, because the data tie the plan to demand the model
cannot see. And asked about a change in the regular price, the model answers
with the response to a temporary cut, about five to six and a half times the
truth. Molecular
attributions in MolSanity~\cite{molsanity} came apart from the truth for a
similar reason: a model can be faithful to what it learned and what it learned
can be wrong.

\paragraph{What this means for a merchant reading a lift}
In the known-answer study, a lift read from a model fitted to the store's own
history inherited its planning habits. If promotions go into weeks that would sell well
anyway, the model's promotion lift overstates what the promotion adds, and an
accurate forecast is no protection: in the known-answer study the correct
world-model seeds and LightGBM forecast better when plans followed demand than
when they were made at random. The check against what sold, the
one most often run in practice, agrees with the overstated lift for the same
reason. Two things help. Agreement between models fitted to the same data
should not be read as confirmation. And a model's answer to a question it was
not trained on, here a lasting change in the regular price from a model that
mostly saw temporary cuts, should be checked against a measured response before
it is used to set prices.

\paragraph{Which checks to run}
The known-answer study ranks the checks. The price-direction check (C4) is cheap
and caught every model blind to price, though it also flags LightGBM, whose
response to price comes in steps. Per-lever agreement with a second model,
an exploratory comparison, caught lifts given to the wrong lever, which every
check on the total lift missed. Construction (C1) and seed stability (C3) are worth confirming once.
The check against what sold (C5) and agreement with LightGBM on the total lift
(C6) confirm agreement, not correctness. Randomisation (C2), the standard sanity
check for attributions~\cite{adebayo2018}, does not test an explanation built
into the model: re-initialising every part left the lifts' ranking largely in
place. A version on the lift as a share of the baseline, or within planned weeks
only, may work; we did not test it. A measured cause (C7) is the only check that
found the regular-price error, though with six spells it could not tell a
missing price response from a correct one. We expect that a measured promotion
effect would be needed to find the bias from planning into strong weeks; we did
not test this, and public retail data offer none.

\paragraph{Ground truth on public data is scarce and noisy}
Dominick's zone price deviations give a pooled
response in line with the one the chain's own experiments reported, but single episodes are noise, intervals
clustered by a handful of stores are far too narrow, and prices start moving
before the detector sees them. A retailer running its own tests would do better
with the design under its control: stores assigned at random, a known start
date, and enough stores per arm for honest intervals.

\paragraph{The audit of the saved versions}
The real versions behaved like the known-answer models in the two ways that
matter most. Checks against the model and against another model were not
shown to track whether the lift is right, and the one comparison with a
measured cause found both models stating a larger price response than the
measured one, as the known-answer world model did, though the measured interval
is wide. The real versions also showed two things the generated data did not.
First, the lift was not zero when the audit planned nothing on most Dominick's
and Frat versions, and we could not fully trace why, so an audit's definition
of nothing planned needs checking against the model's before C1 is read.
Second, the price-direction check failed on most Dominick's and M5 versions
because their forecasts fall after a deeper cut in about half the item-weeks or
more, a fault no accuracy score revealed and one a merchant setting prices from
the model would meet directly. Across versions, passing more checks did not
clearly go with forecasting better. For a retailer this means two separate
tests: an accuracy comparison to choose the forecast, and checks of the
explanation, with a price-direction test and a measured cause among them,
before its lifts are used to plan.

\section{Conclusion}
\label{sec:conclusion}

Plan-lift explanations built into a forecasting model can be faithful, stable
and in agreement with sales and with another model, and still be wrong about
what a plan causes. In generated store data with known effects, the checks that
compare an explanation with the model, the data or another model caught a model
blind to price and, compared lever by lever, lifts given to the wrong lever.
They did not catch lifts inflated by planning into strong weeks, or a model that
answers a question about the regular price with the response to a temporary
cut. Only a measured cause found the second. On public retail data the one such
cause we could measure, Dominick's zone price deviations, gives a pooled response
near the chain's experiments over the full spell, but imprecise.
In the audit of \rVersions{} saved versions on public data, no version passed
every check on M5 or Dominick's, both the world model and LightGBM stated about
twice the measured response to a regular price change (the measured interval is
wide and overlaps both), and the number of checks
a version passed was not clearly related to how well it forecast. On these
models and data sets, an accurate forecast was not evidence of a correct
explanation, and the explanation has to be tested on its own.
